% anexos/anexoB_leer_errores.tex
\chapter{Leer un error de Python}\label{cap:leer_errores}

Cometer errores al programar es inevitable, y no es una mala señal: es parte
del trabajo. Lo que distingue a quien programa con soltura no es que no se
equivoque, sino que sabe leer lo que Python le dice cuando algo falla. Este
anexo enseña a leer esos mensajes con los errores más frecuentes al construir
la calculadora. Todos los mensajes que se muestran son reales: los produjo el
código del libro.

% ---------------------------------------------------------------------------
\section{Cómo se lee un mensaje de error}\label{sec:b_leer}

Supongamos que, al describir una armadura, escribimos por descuido un nodo que
no existe: la barra 2 del ejemplo A va al nodo 4, pero la armadura solo tiene
tres nodos.
\fragmento{b_clave}
Python se detiene y muestra este mensaje:
\salida{b_clave}

El mensaje se llama \termino{traza}{traceback} y se lee \emph{de abajo hacia
arriba}:
\begin{enumerate}
  \item \textbf{La última línea dice qué pasó.} Primero el tipo de error,
    \py{KeyError}, y después el detalle: la clave \py{4} no existe en el
    diccionario.
  \item \textbf{Las líneas anteriores dicen dónde pasó.} Cada bloque que
    empieza con \py{File} es un paso. El primero es nuestra celda; el
    segundo, la línea 10 de la función que llamamos,
    \py{calcular_geometria}, que busca el nodo final de cada barra. Los
    signos \py{~}
    y \py{^} subrayan la parte exacta de la línea que
    falló.
  \item \textbf{El primer bloque es nuestra celda.} Ahí está la línea que
    escribimos nosotros y que desencadenó el error.
\end{enumerate}

Que el error aparezca dentro de una función no significa que la función esté
mal: casi siempre el problema está en los datos que le dimos. Aquí,
\py{calcular_geometria} buscó el nodo final de una barra y no lo encontró,
porque nuestro diccionario \py{barras_mal} usa un nodo inexistente.

\begin{parasabermas}
Colab muestra la misma información con otro formato: agrega colores, marca la
línea que falló con una flecha \py{---->} y nombra las celdas de otra
manera. Se lee igual, de abajo hacia arriba. Colab también puede ofrecer un
botón para que una inteligencia artificial explique el error; es una ayuda
útil, pero conviene aprender a leer el mensaje por nuestra cuenta.
\end{parasabermas}

% ---------------------------------------------------------------------------
\section{Un nombre mal escrito}\label{sec:b_nombre}

El error más frecuente de todos es escribir mal un nombre:
\fragmento{b_nombre}
\salida{b_nombre}
Python no conoce \py{geometira} y, como el nombre se parece a uno que sí
existe, nos sugiere el correcto (\emph{Did you mean}, ¿quisiste decir?).

Un \py{NameError} aparece también cuando usamos una variable o una función
cuya celda no hemos ejecutado: por ejemplo, al abrir de nuevo un notebook,
porque Colab no recuerda nada de la sesión anterior. La solución es ejecutar
las celdas anteriores, o todas con \textsf{Entorno de ejecución > Ejecutar
todas}.

% ---------------------------------------------------------------------------
\section{Errores de escritura}\label{sec:b_escritura}

Algunos errores impiden que Python entienda el código, y aparecen antes de que
se ejecute una sola línea. El más común es olvidar los dos puntos al final de
un \py{for}, un \py{if} o un \py{def}:
\fragmento{b_sintaxis}
\salida{b_sintaxis}
El mensaje no empieza con \py{Traceback}, porque el código no llegó a
ejecutarse: indica la línea y, con el signo \py{^}, el
lugar donde Python esperaba los dos puntos.

El otro error de escritura frecuente es olvidar la sangría del bloque que
sigue a un \py{for}:
\fragmento{b_sangria}
\salida{b_sangria}

% ---------------------------------------------------------------------------
\section{Datos equivocados}\label{sec:b_datos}

Si llamamos a una función con menos argumentos de los que necesita, el
mensaje dice cuál falta. Aquí olvidamos la longitud de la barra:
\fragmento{b_argumentos}
\salida{b_argumentos}

Un \py{IndexError} indica que pedimos una posición que no existe. En la
calculadora casi siempre se debe a confundir la numeración desde 1 con los
índices desde 0. Queremos el desplazamiento vertical del nodo 3, que es el
grado de libertad 6, y lo pedimos con \py{U[6]}:
\fragmento{b_indice}
\salida{b_indice}
\py{U} tiene seis elementos, con índices de 0 a 5. El grado de libertad 6
está en la posición 5, y la forma segura de obtenerla es con la función
\py{indices} del \cref{cap:grados_de_libertad}.

Una división por cero suele indicar dos nodos con las mismas coordenadas: la
barra que los une tiene longitud nula.
\fragmento{b_division}
\salida{b_division}

% ---------------------------------------------------------------------------
\section{Los errores de la calculadora}\label{sec:b_calculadora}

Algunos errores los produce la propia calculadora, a propósito, cuando
detecta datos que no tienen solución. Si a la armadura le falta un apoyo,
$\mK_{ll}$ es singular y \py{resolver_desplazamientos} se detiene con un
mensaje propio (\cref{cap:desplazamientos_reacciones}):
\fragmento{b_inestable}
\salida{b_inestable}
La traza recorre tres pasos: nuestra celda llamó a \py{analizar_armadura},
esta llamó a \py{resolver_desplazamientos}, y ahí se detectó el problema. El
mensaje está en español porque lo escribimos nosotros.

% ---------------------------------------------------------------------------
\section{Los errores que no avisan}\label{sec:b_silenciosos}

Los errores más peligrosos son los que no producen ningún mensaje: una carga
en la unidad equivocada, una coordenada mal escrita o un apoyo en el nodo
equivocado dan resultados, pero resultados falsos. Contra ellos solo sirve la
verificación: comprobar el equilibrio, comparar el dibujo con la armadura y
preguntarse si los números tienen sentido físico, como hicimos a lo largo
del libro. La \cref{tab:errores_frecuentes} del \cref{cap:mi_calculadora}
resume los síntomas más comunes.

\begin{estrategia}
Cuando aparece un error:
\begin{enumerate}
  \item Leer la última línea: el tipo de error y su detalle.
  \item Subir hasta el primer bloque y buscar la línea de nuestra celda que
    lo desencadenó.
  \item Revisar los datos de esa línea: nombres, números de nodo y de barra,
    cantidad de argumentos.
  \item Si el mensaje no basta, imprimir con \py{print} los valores que usa
    la línea, para ver cuál no es el que esperábamos.
  \item Corregir, y ejecutar de nuevo la celda y las que dependen de ella.
\end{enumerate}
\end{estrategia}
